\everymath{\displaystyle}

%%%%%%%%%%%%%%%%%%%%%%%%%%%%%%%%%%%%%%%%
%           main body
%%%%%%%%%%%%%%%%%%%%%%%%%%%%%%%%%%%%%%%%

%-----------------------------------
%	TITLE PAGE
%-----------------------------------

\title[第七章\\
定积分]{\texorpdfstring{\LARGE \bfseries 第七章\quad 定积分}{第七章 定积分}} % The short title appears at the bottom of every slide, the full title is only on the title page
\author[\smaller \S 7.6\\
定积分数值计算] % Your institution as it will appear on the bottom of every slide, maybe shorthand to save space
{\Large \bfseries \S 7.6 定积分的数值计算}
% \author{Singular Value Decomposition} % Your name
% \date{\today} % Date, can be changed to a custom date
\date{}

%------------------------------------------------

\begin{document}

%------------------------------------------------

\begin{frame}

\titlepage % Print the title page as the first slide

\end{frame}

%------------------------------------------------

% \begin{frame}
% \frametitle{内容提要} % Table of contents slide, comment this block out to remove it
% \tableofcontents % Throughout your presentation, if you choose to use \section{} and \subsection{} commands, these will automatically be printed on this slide as an overview of your presentation
% \end{frame}

%-----------------------------------
%	PRESENTATION SLIDES
%-----------------------------------

%------------------------------------------------

\begin{frame}
\frametitle{绪言}

{
\larger[1]
正如之前所提到的, 在实际应用中, 我们很难直接使用定积分的定义来计算定积分的数值. 因此, 我们需要一些数值方法来近似计算定积分的数值. 这只是一些初步的介绍, 后续还会有专门的课程, 例如数值分析课程, 来系统地介绍定积分的数值计算方法.
}

\end{frame}

%------------------------------------------------

\section{Newton-Cotes 求积公式}

%------------------------------------------------

\begin{frame}
\frametitle{待写}

\end{frame}

%------------------------------------------------

\section{复化求积公式}

%------------------------------------------------

\begin{frame}
\frametitle{待写}

\end{frame}

%------------------------------------------------

\section{Gauß 求积公式}

%------------------------------------------------

\begin{frame}
\frametitle{待写}

\end{frame}

%------------------------------------------------

\end{document}
